\section{Blowups of products of projective spaces}\label{sec:mixed}

We first study blowups of products of projective spaces along invariant
codimension-two centres. A simple graph specifies the two factors meeting
each centre. The main result of this section is
Theorem~\ref{mixed:tree}: on a tree, stability against all tangent
subsheaves is equivalent to one divisor-degree inequality at each vertex.
The proof combines a rank formula for the rays with a divergence identity
for the anticanonical polytope. Section~\ref{sec:intrinsic} interprets
the positive terms in that identity as exceptional intersection numbers;
Sections~\ref{sec:mixed-branching}--\ref{sec:mixed-tree-threshold} then
estimate them to obtain uniform stability and local instability theorems.

\subsection{The fan and its primitive collections}

Let $G=(V,E)$ be a finite nonempty simple graph, and write
$k=|V|$, $m=|E|$, and $d_v=\deg_G(v)$. Assign a positive integer $b_v$
to each vertex. In
\[
 N=\bigoplus_{v\in V}\ZZ^{b_v},
\]
let $R_v$ be the $b_v+1$ primitive ray generators of the standard fan of
$\PP^{b_v}$ in its summand. Thus $\sum_{r\in R_v}r=0$. Suppose
$d_v\le b_v+1$, and assign the edges incident to $v$ injectively to rays
$r_{v,e}\in R_v$. For every $e=\{v,w\}$, subdivide the original two-cone
$\langle r_{v,e},r_{w,e}\rangle$ by the ray
\begin{equation}\label{mixed:exceptional-ray}
 s_e=r_{v,e}+r_{w,e}.
\end{equation}
Denote the resulting fan by $\Sigma(G,b)$ and its variety by $Y(G,b)$.
The injectivity condition makes these pairs of original rays disjoint.
All choices of the assignments are related by permutations of the rays
within the projective-space factors and hence give isomorphic fans.

\begin{proposition}\label{mixed:geometry}
The fan $\Sigma(G,b)$ is independent of the order of subdivision and is
smooth and projective. With $n=\sum_v b_v$, one has
\begin{equation}\label{mixed:dimension}
 \dim Y(G,b)=n,\qquad \rho(Y(G,b))=k+m.
\end{equation}
Its primitive collections are precisely
\begin{align}
 \mathcal C_e&=\{r_{v,e},r_{w,e}\},\qquad e=\{v,w\},
 \label{mixed:edge-collection}\\
 \mathcal P_v(J)&=
 \bigl(R_v\setminus\{r_{v,e}:e\in J\}\bigr)
 \cup\{s_e:e\in J\},\qquad J\subseteq E(v),
 \label{mixed:vertex-collection}
\end{align}
where $E(v)$ is the set of edges incident to $v$. Their primitive relations
are
\begin{equation}\label{mixed:primitive-relations}
 r_{v,e}+r_{w,e}=s_e,\qquad
 \sum_{r\in\mathcal P_v(J)}r
   =\sum_{\substack{e=\{v,w\}\in J}}r_{w,e}.
\end{equation}
Every $Y(G,b)$ is weak Fano, and
\begin{equation}\label{mixed:fano}
 Y(G,b)\text{ is Fano}
 \quad\Longleftrightarrow\quad b_v\ge d_v\quad(v\in V).
\end{equation}
If $G$ is connected, $Y(G,b)$ admits no torically locally trivial bundle
presentation with positive-dimensional complete fibre and
positive-dimensional base.
\end{proposition}

\begin{proof}
In an original maximal cone, the subdivision centres that it contains
use disjoint pairs of basis vectors. The subdivisions are therefore
products of subdivisions of two-dimensional orthants. They commute and
remain unimodular. Each remaining original centre persists after the
other subdivisions. Every step is a toric blowup along a smooth invariant
centre, so the resulting fan is projective. It has $n+k+m$ rays, proving
\eqref{mixed:dimension}.

We first describe all its faces. For a subset $A$ of the final rays,
write $B$ for its original rays and $J$ for the edges whose exceptional
rays belong to $A$. Then $A$ spans a cone if and only if
\begin{enumerate}
\item $B$ contains no endpoint pair $\mathcal C_e$;
\item the union of $B$ with both endpoints of every $e\in J$ contains
      no complete set $R_v$.
\end{enumerate}
Indeed, the smallest original cone containing $s_e$ contains both its
endpoints, whereas the original endpoint pair is removed by subdivision.
This proves necessity. For sufficiency, the second condition places the
expanded set in an original product cone. Within that cone, the
subdivisions act on disjoint coordinate pairs, and the first condition
is exactly the condition for the chosen rays to lie in one of the
resulting cones.

The sets in \eqref{mixed:edge-collection} are minimal nonfaces. The
expansion of $\mathcal P_v(J)$ contains $R_v$, and deleting any one member
removes its unique representative of one ray of $R_v$. At any other
vertex, that expansion contains at most one ray, because $G$ is simple;
it cannot contain the full set of at least two rays of that factor.
There is also no original endpoint pair in $\mathcal P_v(J)$. The face
criterion proves its minimality, including when every member is an
exceptional ray. Conversely, a nonface containing no $\mathcal C_e$
expands to a set containing some $R_v$. Choosing one original or
exceptional representative for each ray of $R_v$ produces a subset
$\mathcal P_v(J)$. The list is complete.

The vector identities \eqref{mixed:primitive-relations} follow from
$\sum_{r\in R_v}r=0$. On the right of the second identity there is at
most one ray in each neighbouring factor. No two of these rays form
another centre, since incidence rays are assigned injectively. They
therefore span a cone of the final fan, and their positive sum lies in
its relative interior. These are the primitive relations, of degrees
$1$ and $b_v+1-|J|$, respectively. Batyrev's primitive-collection
criterion for nefness and ampleness
\cite[Theorem~6.4.9]{CoxLittleSchenck2011} proves that $-K$ is nef and
gives \eqref{mixed:fano}. The anticanonical polytope contains the origin
in its interior and is bounded by completeness, so $-K$ is also big.

Finally, the cone complex of a torically locally trivial bundle is the
join of its fibre and base cone complexes. Its rays partition into fibre
rays and one lift of each base ray, and every primitive collection lies
entirely in one part. This follows on the invariant affine charts, where
the fan is the product of the fibre fan and the base cone; the
identifications agree on common faces. In our fan, the primitive
collection $R_v=\mathcal P_v(\varnothing)$ joins all original rays of
one factor. The collections $\mathcal C_e$ join neighbouring factors,
and $\mathcal P_v(\{e\})$ joins $s_e$ to $b_v\ge1$ original rays.
For connected $G$, all rays consequently lie in one equivalence class
generated by membership in a common primitive collection. They admit
no partition into two nonempty parts respecting every primitive
collection. This excludes every bundle presentation of the stated kind,
including alternative choices of fibre and base.
\end{proof}

The last assertion concerns locally trivial toric bundles; it does not
exclude toric morphisms with special fibres or arbitrary non-equivariant
fibrations. Simplicity of $G$ and injectivity of the incidence assignments
are part of the construction. For example, two parallel edges using the
opposite pairs of rays of $\PP^1\times\PP^1$ give the degree-six del Pezzo
surface, which does not satisfy \eqref{mixed:fano} with these vertex
dimensions.

\subsection{Divisor degrees and a divergence identity}

For the remainder of this section assume $b_v\ge d_v$, and put
$Y=Y(G,b)$. Use the normalized degrees of Section~\ref{sec:stability}:
\begin{equation}\label{mixed:normalized-degrees}
 \delta_r=\frac{nD_r\cdot(-K_Y)^{n-1}}{(-K_Y)^n},
 \qquad \sum_r\delta_r=n.
\end{equation}
They are all positive. Set $C_e=\delta_{s_e}$. There is a common value
$A_v$ for the degrees of the unassigned original rays at $v$, and
\begin{equation}\label{mixed:degree-relations}
 \delta_{r_{v,e}}=A_v-C_e,\qquad
 \sum_v(b_v+1)A_v-\sum_eC_e=n.
\end{equation}
To see this, project the relation $\sum_r\delta_r r=0$ of
Lemma~\ref{lem:degreerelation} onto the $v$-th lattice summand. The
coefficient of an assigned ray is $\delta_{r_{v,e}}+C_e$, and that
of an unassigned ray is its degree. The unique relation among $R_v$
has all coefficients equal, proving the first assertion. Summing all
ray degrees gives the second. In particular,
\begin{equation}\label{mixed:degree-positive}
 0<C_e<A_v\qquad(e\in E(v)).
\end{equation}

There are $q_v=b_v+1-d_v\ge1$ unassigned rays at $v$. Choose one as
the negative sum ray and write the others as an integral basis
$r_{v,1},\ldots,r_{v,b_v}$. If $M=\operatorname{Hom}(N,\ZZ)$, the affine
coordinates $a_{v,j}=1+\langle x,r_{v,j}\rangle$ on $M_\RR$ preserve
lattice measure. We write $a_{v,e}$ for the coordinate assigned to the
incidence $(v,e)$. The anticanonical polytope $P$ is
\begin{equation}\label{mixed:polytope}
 a_{v,j}\ge0,\qquad \sum_{j=1}^{b_v}a_{v,j}\le b_v+1,
 \qquad a_{v,e}+a_{w,e}\ge1\quad(e=\{v,w\}).
\end{equation}
Throughout, volumes are lattice volumes without factorial normalization.
If $F_r$ is the facet corresponding to $r$, then
\begin{equation}\label{mixed:facet-ratio}
 (-K_Y)^n=n!\operatorname{vol}(P),\qquad
 \delta_r=\frac{\operatorname{vol}(F_r)}{\operatorname{vol}(P)}.
\end{equation}
The second formula follows from
$D_r\cdot(-K_Y)^{n-1}=(n-1)!\operatorname{vol}(F_r)$.

\begin{lemma}\label{mixed:flux}
For an edge $e=\{v,w\}$, define
\[
 J_{v,e}=\frac{1}{\operatorname{vol}(P)}
     \int_{F_{s_e}}a_{v,e}\,d\operatorname{vol}.
\]
Then
\begin{equation}\label{mixed:flux-identity}
 (b_v+1)A_v=b_v+\sum_{e\in E(v)}J_{v,e},\qquad
 J_{v,e}+J_{w,e}=C_e,
\end{equation}
and $0<J_{v,e}<C_e<A_v$. In particular, $A_v<1$ at every vertex
of degree one, for every positive value of $b_v$.
\end{lemma}

\begin{proof}
Apply the divergence theorem to the vector field equal to $a_v$ in the
$v$-th coordinate block and zero elsewhere. Its divergence is $b_v$.
The coordinate facets have zero flux. The facet
$\sum_j a_{v,j}=b_v+1$ has flux $(b_v+1)$ times its lattice volume,
and the exceptional facet $F_{s_e}$ has flux
$-\int_{F_{s_e}}a_{v,e}\,d\operatorname{vol}$. All remaining facets have
zero flux. For a primitive integral normal, its Euclidean length in
hypersurface measure cancels the reciprocal factor in the unit-normal
flux. Division by $\operatorname{vol}(P)$ proves the first identity.

On $F_{s_e}$ one has $a_{v,e}+a_{w,e}=1$. Its relative interior has
both coordinates strictly between zero and one, since intersection
with either coordinate facet has smaller dimension. This proves the
second identity and the strict bounds on $J_{v,e}$; the remaining bound
is \eqref{mixed:degree-positive}. At a leaf, the first identity gives
$(b_v+1)A_v=b_v+J_{v,e}<b_v+A_v$, hence $A_v<1$.
\end{proof}

\subsection{A cographic rank formula}

Construct a multigraph $H$ with vertices $o$, the vertices $v$ of $G$,
and one vertex $z_e$ for each edge of $G$. Its edges are indexed by the
rays of $\Sigma(G,b)$ as follows:
\begin{itemize}
\item an assigned ray $r_{v,e}$ corresponds to an edge $z_e\to v$;
\item each unassigned original ray at $v$ corresponds to an edge $o\to v$;
\item the exceptional ray $s_e$ corresponds to an edge $o\to z_e$.
\end{itemize}
The orientations specify an incidence matrix and do not affect graph
connectivity. Thus $H$ has $k+m+1$ vertices and $n+k+m$ edges. It is
connected, even if $G$ is disconnected: each $z_e$ joins the root, each
nonisolated $v$ joins some $z_e$, and an isolated $v$ has unassigned
rays joining it to the root.

\begin{proposition}\label{mixed:rank}
For a subset $S$ of the rays, let $H\setminus S$ denote the graph obtained
by deleting its corresponding edges and retaining all vertices. Then
\begin{equation}\label{mixed:rank-formula}
 r(S)=|S|+1-c(H\setminus S),
\end{equation}
where $c$ counts connected components, including isolated vertices.
Consequently the ray matroid is the cographic matroid of $H$: its bases
are exactly the complements of spanning trees of $H$.
\end{proposition}

\begin{proof}
Let $B$ be the signed incidence matrix of $H$, with positive entries at
the heads of oriented edges, and let $R:\QQ^{E(H)}\to N_\QQ$ send each
edge basis vector to its ray. The row at $v$ maps to
$\sum_{r\in R_v}r=0$, and the row at $z_e$ maps to
$s_e-r_{v,e}-r_{w,e}=0$. Thus the row space of $B$ lies in $\ker R$.
Both spaces have dimension $k+m$: the incidence matrix has that rank
because $H$ is connected, and the original rays span $N_\QQ$. Hence
$\ker R=\operatorname{im}B^{\mathsf T}$.

A relation supported on $S$ is therefore $B^{\mathsf T}f$, where $f$
is constant on every component of $H\setminus S$. Modulo the global
constants, these functions have dimension $c(H\setminus S)-1$.
Subtracting this relation dimension from $|S|$ proves
\eqref{mixed:rank-formula}. Rational and complex ranks agree. A subset
of $n$ rays is a basis precisely when its complement is connected;
that complement has $k+m$ edges on $k+m+1$ vertices, so it is a
spanning tree.
\end{proof}

\subsection{The stability criterion for trees}

\begin{theorem}\label{mixed:tree}
Suppose $G$ is a tree with at least one edge and $b_v\ge d_v$ for all
$v$. Then $T_{Y(G,b)}$ is anticanonically stable if and only if
$A_v<1$ for every vertex, and is semistable if and only if $A_v\le1$
for every vertex. Thus it is strictly semistable when $\max_v A_v=1$
and unstable when $\max_v A_v>1$.
\end{theorem}

\begin{proof}
Let $K$ be the tree obtained by subdividing every edge of $G$ at $z_e$,
so $H$ is obtained from $K$ by adjoining the root and its edges. Give
each edge of $H$ its normalized ray degree. There are $q_v$ root edges
of weight $A_v$ at $v$, one root edge of weight $C_e$ at $z_e$, and
the edge $v z_e$ has weight $A_v-C_e$.

For a ray subset $S$, set $E'=E(H)\setminus S$, write $\delta(E')$
for its total weight, and put
\[
 \beta(E')=|E'|-|V(H)|+c(E').
\]
This is the cycle rank, with all vertices retained. By
Proposition~\ref{mixed:rank}, $r(S)=n-\beta(E')$. Since the total
weight is $n$, the slope inequality becomes
\begin{equation}\label{mixed:complement-test}
 \delta(E')>\beta(E')\qquad\text{when }0<r(S)<n,
\end{equation}
with a weak inequality for semistability.

For a nonempty connected vertex set $Q$ of $K$, let $q(Q)$ count all
root edges at its vertices, including multiplicity, and let $W(Q)$ be
the sum of their weights and the weights of all $K$-edges internal to
$Q$. Define
\[
 D(Q)=W(Q)-(q(Q)-1).
\]
For $Q=K$ this is zero, because the cycle rank and total weight of $H$
both equal $n$. For $Q=\{z_e\}$ it is $C_e>0$.

Otherwise the original vertices in $Q$ form a nonempty connected subtree
$I$ of $G$. All exceptional vertices on internal edges of $I$ belong
to $Q$. Let $\partial I$ be the edges with exactly one endpoint in $I$,
write $v(e)$ for that endpoint, and set $t_e=1$ if $z_e\in Q$, and
$t_e=0$ otherwise. Counting the original-ray contributions and then
using \eqref{mixed:flux-identity} gives
\begin{align*}
 W(Q)
 &=\sum_{v\in I}(b_v+1)A_v
   -\sum_{e\in\partial I}(1-t_e)A_{v(e)}
   -\sum_{e\text{ internal to }I}C_e\\
 &=\sum_{v\in I}b_v+\sum_{e\in\partial I}J_{v(e),e}
   -\sum_{e\in\partial I}(1-t_e)A_{v(e)}.
\end{align*}
There are $|I|-1$ internal edges, and
$\sum_{v\in I}d_v=2(|I|-1)+|\partial I|$. Consequently
\[
 q(Q)-1=\sum_{v\in I}b_v-|\partial I|+
                 \sum_{e\in\partial I}t_e,
\]
and hence
\begin{equation}\label{mixed:boundary-deficit}
 D(Q)=\sum_{e\in\partial I}
 \bigl[J_{v(e),e}+(1-t_e)(1-A_{v(e)})\bigr].
\end{equation}
Assume now $A_v\le1$ for every $v$. If $I$ is proper, its boundary is
nonempty, so this sum is strictly positive. If $I$ contains every original
vertex, connectivity forces $Q=K$. Together with the exceptional singleton
case, this proves $D(Q)>0$ for every nonempty proper connected $Q$.

The selected non-root edges of $E'$ form a forest in $K$. Each component
has a connected vertex set $Q$ and contains every $K$-edge internal to
that set, since $K$ is a tree. If it has $h$ selected root edges, counted
with multiplicity, it contributes $\max(h-1,0)$ to $\beta(E')$. This
remains true when several components join the common root: each component
with a root edge adds one connection and its remaining root edges add
cycles. For $h=0$, its contribution to weight minus cycle rank is
nonnegative, and positive if it contains an edge. For $h\ge1$, completing
the missing root edges gives the exact contribution
\begin{equation}\label{mixed:root-completion}
 D(Q)+\sum_{r\text{ a missing root edge at }Q}(1-\delta_r).
\end{equation}
All root degrees are at most one; indeed $C_e<A_v\le1$. Thus every
component contribution is nonnegative, proving the weak inequalities
in \eqref{mixed:complement-test}.

If every $A_v<1$, all root degrees are strictly below one. Contributions
from proper components with root edges are already positive. A full
component $Q=K$ is also positive if a root edge is missing, by
\eqref{mixed:root-completion}; if none is missing, $E'=E(H)$ and
$S=\varnothing$. Components without root edges contribute positively
whenever they contain an edge. The only remaining graph with no positive
contribution is the empty $E'$, which has $r(S)=n$. Both exceptions lie
outside the tested ranks. This proves every strict inequality.

Conversely, an unassigned ray of degree $A_v$ belongs to a rank-one
subspace of weight at least $A_v$. This is a proper subspace because
$n\ge2$. Stability therefore requires $A_v<1$, and semistability
requires $A_v\le1$. These necessity statements, the inequalities just
proved, and Lemma~\ref{lem:klyachko} give all three alternatives.
\end{proof}

The reduction to vertex inequalities uses two different features of the
construction. Proposition~\ref{mixed:rank} describes all tangent subspace
ranks for any simple graph. For trees, the boundary formula
\eqref{mixed:boundary-deficit} expresses the remaining inequalities as
positive sums of vertex and exceptional contributions. In particular,
the scalar criterion is not obtained by disregarding higher-rank
subsheaves. The tree hypothesis is used in this passage from local
degrees to all subspace inequalities.

\begin{corollary}\label{mixed:single-edge}
For every $p,q\ge1$, the blowup of $\PP^p\times\PP^q$ along the invariant
codimension-two subvariety corresponding to a cone generated by one ray
from each factor has anticanonically stable tangent bundle.
\end{corollary}

\begin{proof}
Both vertices of its graph have degree one. Lemma~\ref{mixed:flux}
gives both strict inequalities required by Theorem~\ref{mixed:tree}.
\end{proof}

For $(p,q)=(1,3)$, this recovers
\cite[Proposition~5.6.1(2)]{DasguptaDeyKhan2020}.

\subsection{Stable paths and the largest Picard number}

\begin{theorem}\label{mixed:paths}
Let $G$ be a path with $L\ge1$ edges, endpoint dimensions $p,q\ge1$,
and dimension two at every internal vertex. Then $T_{Y(G,b)}$ is
anticanonically stable, and
\begin{equation}\label{mixed:path-dimension}
 \dim Y(G,b)=2L+p+q-2,\qquad \rho(Y(G,b))=2L+1.
\end{equation}
At each internal vertex its unassigned normalized degree satisfies
$A_v\le24/25$.
\end{theorem}

\begin{proof}
Fix the incidence coordinate $a$ on one side of an edge and integrate
the coordinates of the branch on the other side. Denote its volume by
$\phi(a)$. We shall use the following properties:
\begin{equation}\label{mixed:incoming-bound}
 \phi(a)>0,\qquad \phi(a)=c:=\phi(1)\ (a\ge1),\qquad
 \phi(a)\ge ac\ (0\le a\le1).
\end{equation}
A leaf of dimension $b$ gives
\[
 \phi_b(a)=\frac{(b+\min(a,1))^b}{b!}.
\]
Indeed, translating its assigned coordinate by the lower bound
$\max(0,1-a)$ leaves a simplex of that capacity. Bernoulli's inequality
gives
\[
 \frac{\phi_b(a)}{\phi_b(1)}
 =\left(1-\frac{1-a}{b+1}\right)^b
 \ge1-\frac{b}{b+1}(1-a)\ge a,
\]
so \eqref{mixed:incoming-bound} holds for every endpoint dimension.

At a degree-two vertex, integration of the coordinate assigned to the
parent transforms an incoming function by
\begin{equation}\label{mixed:path-operator}
 g(a)=\int_0^{2+\min(a,1)}
             (2+\min(a,1)-t)\phi(t)\,dt.
\end{equation}
For a positive input constant beyond one, set
$I_0=\int_0^1\phi(t)\,dt$ and $I_1=\int_0^1t\phi(t)\,dt$.
On $0\le a\le1$, the output is $g(a)=g_0+g_1a+g_2a^2$, where
\[
 g_0=2I_0-I_1+c/2,\qquad g_1=I_0+c,\qquad g_2=c/2.
\]
All coefficients are positive, and $g_0-g_2=2I_0-I_1>0$. Hence
\[
 g(a)-ag(1)=(1-a)(g_0-ag_2)\ge0.
\]
The output is constant beyond one. Induction along each branch proves
\eqref{mixed:incoming-bound} for every incoming function, including
branches whose first input is not quadratic.

At an internal vertex let $\phi$ and $\psi$ be the two incoming
functions. The branches have disjoint coordinate blocks, so their
product is the integrand over $a,b\ge0$, $a+b\le3$. For either input
use subscripts to distinguish $c$, $I_0$ and $I_1$. Splitting this
simplex at $a=1$ and $b=1$ gives its full volume $V$ and its unassigned
facet volume $F$ as
\begin{align}
 V={}&I_{0,\phi}I_{0,\psi}
   +c_\psi(2I_{0,\phi}-I_{1,\phi})
   +c_\phi(2I_{0,\psi}-I_{1,\psi})
   +\tfrac12c_\phi c_\psi,\label{mixed:path-volume}\\
 F={}&c_\phi c_\psi+c_\psi I_{0,\phi}+c_\phi I_{0,\psi}.
 \label{mixed:path-facet}
\end{align}
The first term of $V$ is the unit square, the next two are the regions
with just one coordinate at least one, and the last is the remaining
triangle. Eliminating one coordinate from $a+b=3$ gives the lattice
measure in the formula for $F$.

Define
\[
 i_\phi=I_{0,\phi}/c_\phi,\qquad
 h_\phi=(I_{0,\phi}-I_{1,\phi})/c_\phi,
\]
and similarly for $\psi$. From \eqref{mixed:incoming-bound},
$i_\phi\ge1/2$ and $h_\phi\ge1/6$. Equations
\eqref{mixed:path-volume}--\eqref{mixed:path-facet} become
\[
 \frac{F}{c_\phi c_\psi}=1+i_\phi+i_\psi,\qquad
 \frac{V-F}{c_\phi c_\psi}
   =i_\phi i_\psi+h_\phi+h_\psi-\frac12.
\]
Write $i_\phi=1/2+x$, $i_\psi=1/2+y$, $h_\phi=1/6+u$ and
$h_\psi=1/6+v$, with $x,y,u,v\ge0$. Then
\[
 \frac{24V-25F}{c_\phi c_\psi}
   =24xy+11(x+y)+24(u+v)\ge0.
\]
Thus \eqref{mixed:facet-ratio} gives $A_v=F/V\le24/25<1$ at every
internal vertex. Lemma~\ref{mixed:flux} gives $A_v<1$ at both endpoints.
Theorem~\ref{mixed:tree} proves stability, also for $L=1$, when there
is no internal vertex. The dimension and Picard number follow from
Proposition~\ref{mixed:geometry}.
\end{proof}

\begin{corollary}\label{mixed:maximum}
For every $n\ge1$, the largest Picard number among the smooth Fano
$n$-folds $Y(G,b)$ with stable tangent bundle is
\begin{equation}\label{mixed:maximum-value}
 2\left\lfloor\frac n2\right\rfloor+1.
\end{equation}
Here $G$ ranges over the finite nonempty simple graphs with injective
incidence-ray assignments in the construction above.
\end{corollary}

\begin{proof}
If $G$ is disconnected, its fan is the product of the fans for its
connected components. The tangent bundle is a direct sum of at least
two positive-rank bundles; one summand has slope at least the slope of
the whole bundle. Stability therefore forces connectedness. Consequently
$k\le m+1$, whereas the Fano condition gives
\[
 n=\sum_v b_v\ge\sum_v d_v=2m.
\]
It follows that $\rho=k+m\le2m+1\le2\lfloor n/2\rfloor+1$.

For $n=2L\ge2$, take the length-$L$ path with endpoint dimensions
$(1,1)$. For $n=2L+1\ge3$, take endpoint dimensions $(1,2)$ instead.
Theorem~\ref{mixed:paths} gives stability, and
\eqref{mixed:path-dimension} gives equality in \eqref{mixed:maximum-value}.
For $n=1$, the single-vertex graph with dimension one gives $\PP^1$,
whose rank-one tangent bundle is stable and whose Picard number is one.
\end{proof}

This maximum concerns the specified blowups of products of projective
spaces. It is separate from the bundle-class maximum in
Section~\ref{sec:interactions} and does not give an upper bound for
arbitrary smooth toric Fano varieties.
